% Source: Weinberg GR, physical PDF p. 156
%         (printed p. 131).
% Coverage: Chapter 6 epigraph and opening overview.
% Figures/tables/footnotes: none.
% Status: source-reviewed and compile-clean.
% Uncertainties: none.

\begin{flushright}
\begin{minipage}{0.42\textwidth}
\itshape
``When I behold, upon the\\
night's starr'd face, huge\\
cloudy symbols \ldots''

\medskip
\upshape
John Keats, \textit{When I Have\\
Fears That I May Cease\\
to Be}
\end{minipage}
\end{flushright}

We are going to work out the gravitational field equations by applying the
Principle of Equivalence to gravitation itself. As in the last chapter, it is
most convenient to apply this principle by looking for field equations that
are generally covariant and that reduce to the proper form for weak fields.
Thus we must address ourselves to the question: What tensors can be formed
from the metric tensor and its derivatives? In this chapter we treat this as a
purely mathematical problem, as in fact it was treated by Gauss and Riemann;
the information we compile here will then be used in the next chapter to guide
us in our search for the field equations of gravitation.
